\textbf{Задача 2.} Докажите, что если $A\in\p$, то $A^*\in\p$, и что если $A\in\np$, то $A^*\in\np$.

\begin{proof} Пусть $A\in\p$. Пусть на вход дают слово $w\ : \ |w|=n$. Посчитаем динамику 
\[dp[k] = I\left(w[0,k]\in A^*\right)\quad k=0,\ldots,n-1\]

База $dp[0]=I\left(w[0]\in A\right)$.\\
Далее для подсчета $dp[k]$ делаем цикл по $m=1\ldots k -1$ и если найдется $m$, что $dp[m]=1$ и $w[m+1,k]\in A$, то ставим $dp[k] = 1$, иначе ставим $dp[k] = 0$.

Итоговый ответ $dp[n -1]$.

Пусть $A\in\np$, следовательно, существует машина $V(x,y)$ работающая за полиномиальное от входа время, причём
\[\forall x \in \{0,1\}^*\quad x\in A\iff \exists y\ |y|\le p(|x|)\ V(x,y)=1\]

Построим $W(x,y)$ работающую за полиномиальное от входа время, так чтобы
\[\forall x \in \{0,1\}^*\quad x\in A^*\iff \exists y\ |y|\le p(|x|)\ W(x,y)=1\]

Для этого в качестве сертификата $y$ будем давать границы по которым надо разбивать и сертификаты (для машины $V$) для каждого кусочка.

Так как кусочков линейное число, сертификат будет полиномиальной от входа длинны.

Машина $W$ будет просто проверять, что разбиение корректно.
\end{proof}